\section{Post-Freeze Cross-Schema Robustness Evaluation}
\label{sec:cross_schema}
\begin{table*}[t]
\centering
\caption{Cross-Schema Robustness Evaluation across Four Heterogeneous Enterprise-Style Domains ($N=32$ Tasks total, 8 per Schema)}
\label{tab:cross_schema}
\resizebox{\textwidth}{!}{%
\begin{tabular}{llccccccc}
\toprule
\textbf{Enterprise-Style Tool Schema} & \textbf{Domain / Key Constraints} & \textbf{Init $S_0$} & \textbf{Full $S$} & \textbf{Token Savings} & \textbf{1st-Pass (Exp/Amb)} & \textbf{Final Struct.} & \textbf{Semantic} & \textbf{Ambiguous Recovery} \\
\midrule
\texttt{create\_purchase\_order} & Procurement (Regex, enums, nested items) & 339 & 552 & \textbf{-21.3\%} & 50.0\% (4/4, 0/4) & \textbf{100.0\%} & 87.5\% & \textbf{100.0\%} (4/4) \\
\texttt{inventory\_adjustment} & Warehouse Stock (Signed deltas, reason codes) & 299 & 541 & \textbf{-22.4\%} & 50.0\% (4/4, 0/4) & \textbf{100.0\%} & 87.5\% & \textbf{100.0\%} (4/4) \\
\texttt{shipment\_creation} & Freight Logistics (Nested packages, tracking regex) & 337 & 646 & \textbf{-25.2\%} & 50.0\% (4/4, 0/4) & \textbf{100.0\%} & 100.0\% & \textbf{100.0\%} (4/4) \\
\texttt{customer\_refund} & Financial Accounting (Bounds, currency enums) & 292 & 474 & \textbf{-20.5\%} & 50.0\% (4/4, 0/4) & \textbf{100.0\%} & 100.0\% & \textbf{100.0\%} (4/4) \\
\midrule
\textbf{Macro Mean / Aggregate} & \textbf{Cross-Domain Summary ($N=32$)} & \textbf{316.8} & \textbf{553.2} & \textbf{-22.5\%} & \textbf{50.0\%} (16/16, 0/16) & \textbf{100.0\%} & \textbf{93.8\%} & \textbf{100.0\%} (16/16) \\
\bottomrule
\end{tabular}%
}
\end{table*}


To assess whether the LASR-MCP mechanism transfers beyond purchase-order semantics and obtain preliminary evidence of cross-schema transfer, we conducted a post-freeze robustness study across three new synthetic enterprise-style schemas plus a replication benchmark on the primary purchase-order schema (four schemas total) conforming to JSON Schema Draft 2020-12~\cite{jsonschema202012}: (1)~\textbf{\texttt{create\_purchase\_order}} (Procurement replication: 552 full / 339 $C_I^*$ tokens; regexes, enums, budget bounds, nested items); (2)~\textbf{\texttt{inventory\_adjustment}} (Warehouse: 541 / 299 tokens; signed stock deltas, facility codes, reason codes); (3)~\textbf{\texttt{shipment\_creation}} (Logistics: 646 / 337 tokens; multi-carrier enums, tracking regexes, nested packages, weight bounds); and (4)~\textbf{\texttt{customer\_refund}} (Finance: 474 / 292 tokens; refund ceilings, multi-currency enums, return classifications).

\subsection{Evaluation Methodology and Scrutiny Controls}
For each schema, we constructed 8 pre-specified evaluation tasks ($N=32$ total tasks), partitioned into two operational regimes: \textbf{Explicit Tasks} (16 tasks, 4 per schema), where intents furnish required enum tokens and parameters; and \textbf{Ambiguous Tasks} (16 tasks, 4 per schema), intentionally constructed as a failure-induction stress test without explicit cues, starting with lean $C_I^*$ to induce initial violations and test closed-loop reactive recomposition.

To ensure empirical validity, we adhere to strict methodological controls: (1) all 32 tasks, target schemas, and invariant foundations were pre-specified before any model execution or outcome inspection; (2) no task was added, removed, or cherry-picked post-hoc; (3) no hyperparameter or predictor weight was tuned on these new domains; and (4) the original frozen benchmark ($N=32$, commit \texttt{9807340}, SHA-256 \texttt{78c2c88e}) was kept strictly quarantined and untouched.

\subsection{Scope of Transfer: Architecture vs. Predictor}
We delineate transfer scope: \textbf{1) Architecture Transfer}: Closed-loop reactive recovery---validator error interception, keyword extraction, and in-tree recomposition---operates domain-agnostically across schemas. \textbf{2) Schema Foundations}: Each schema instantiates its own invariant floor ($C_I^* \in \{339, 299, 337, 292\}$ tokens) and constraint dictionary $\mathcal{S}_{\text{specs}}$, fixed \textit{a priori}. \textbf{3) Predictor Status}: We evaluate recovery mechanism transfer without retraining/retuning the primary predictor or modifying the frozen benchmark. Explicit tasks test intent-conditioned rules; ambiguous tasks test zero-shot recovery from baseline $C_I^*$. Thus, this experiment evaluates \textbf{architecture and mechanism transfer}, rather than universal zero-shot predictor transfer across arbitrary unseen APIs.

\subsection{Empirical Robustness Results}
Table~\ref{tab:cross_schema} summarizes system performance across all four enterprise-style domains:

\textit{Domain Token Savings}: Across the three new synthetic domains, LASR consistently reduces cumulative token consumption: saving \textbf{22.4\%} on \texttt{inventory\_adjustment}, \textbf{25.2\%} on \texttt{shipment\_creation}, and \textbf{20.5\%} on \texttt{customer\_refund} (macro mean: \textbf{22.7\%}), alongside \textbf{21.3\%} on procurement replication (unrounded macro mean: 22.4\%; aggregate workload reduction: \textbf{-22.5\%}). On multi-turn ambiguous tasks requiring reactive disambiguation, a compact foundation ($C_I \approx 290$--$340$ tok) plus targeted reactive delta ($\Delta S \approx 40$ tok) consumes substantially fewer tokens than transmitting monolithic schemas across every turn.

\textit{Complete Ambiguous Task Recovery (16/16)}: The 50.0\% first-pass rate (16/16 explicit, 0/16 ambiguous) reflects the 50/50 stress-test design. On all 16 ambiguous tasks where first-pass rejections occurred from unconstrained parameters, the multi-error extractor intercepted diagnostics and recomposed missing subtrees into the AST, achieving a \textbf{100.0\% structural recovery rate} on Attempt~2 and \textbf{100.0\% final structural compliance} across all 32 tasks (representing observed stress-test performance, not an unconstrained universal guarantee).

\textit{High Semantic Fidelity}: Evaluated via independent deterministic domain validators checking operational fidelity (routing, quantities, carrier assignment, currency), LASR attained \textbf{87.5\%} to \textbf{100.0\%} semantic correctness (macro aggregate: \textbf{93.8\%}), verifying that compact schema exposure preserves domain execution fidelity.

These findings provide preliminary evidence of cross-schema transfer: LASR's core architectural mechanics transfer effectively across diverse enterprise-style JSON Schema contracts without retuning the recovery algorithm or the primary predictor; schema-specific constraint vocabularies and invariant foundations were instantiated for each schema.
